\section{The workload: foundations of neural inference}
\label{sec:ml}

Every measurement in this survey is of the same computation: a forward pass
through a trained neural network in which the input, and in our case also the
weights, are hidden from the party performing it. A reader cannot judge those
measurements without knowing precisely which function is being computed, which of
its parts are expensive under which cryptographic paradigm, and what an accuracy
figure for it means. This section supplies that. It is deliberately elementary in
the machine-learning direction and precise in the arithmetic direction, because
the arithmetic is what the cryptography charges for.

\subsection{Classification and the inference function}
\label{sec:ml:class}

\begin{definition}
\label{def:classifier}
Let $\mathcal{X} \subseteq \Rr^{n_0}$ be an input space and
$\mathcal{Y} = \{0,\dots,C-1\}$ a finite label set. A \ib{classifier} is a map
$h : \mathcal{X} \to \mathcal{Y}$. It is usually presented as
$h = \argmax \circ\, g$ where $g : \mathcal{X} \to \Rr^{C}$ produces a vector of
\ib{logits} and $\argmax$ returns the index of the largest.
\end{definition}

Training chooses $g$ from data; \ib{inference} evaluates it on one input. This
survey measures inference only. The distinction matters because the two have
different arithmetic: training needs gradients, and therefore division and
accumulation across a batch, whereas inference needs a fixed sequence of affine
maps and nonlinearities. Our artifact does implement private training, and we
report it, but the comparisons in Section~\ref{sec:cross} are of inference.

\begin{definition}
\label{def:ffn}
A \ib{feedforward network} of depth $L$ is a composition
\[
  g \;=\; A_L \circ \sigma \circ A_{L-1} \circ \cdots \circ \sigma \circ A_1 ,
\]
where each $A_\ell : \Rr^{n_{\ell-1}} \to \Rr^{n_\ell}$ is affine,
$A_\ell(\mathbf{u}) = \mathbf{W}_\ell \mathbf{u} + \mathbf{b}_\ell$, and $\sigma$
is a fixed nonlinearity applied coordinatewise. We call $n_\ell$ the \ib{width}
of layer $\ell$ and $L$ the \ib{depth}. The parameters are
$\theta = (\mathbf{W}_\ell, \mathbf{b}_\ell)_{\ell \le L}$.
\end{definition}

Definition~\ref{def:ffn} already exposes the structure that governs every cost in
this survey. The affine maps are \emph{linear} in the shared or encrypted value,
so under any linearly homomorphic representation---an additive secret sharing, or
a lattice ciphertext---they require no interaction and no ciphertext
multiplication beyond scaling by public constants. The nonlinearity $\sigma$ is
not linear, and it is the only place where a protocol must do something
irreducibly expensive. The entire comparison between paradigms is therefore a
comparison of how each one implements $\sigma$, and how many times it must.

\begin{proposition}
\label{prop:nonlin-count}
A network as in Definition~\ref{def:ffn} performs
$\sum_{\ell} n_{\ell-1} n_\ell$ multiply--accumulate operations in its affine maps
and $\sum_{\ell < L} n_\ell$ applications of $\sigma$. The first quantity is
quadratic in width and the second is linear; consequently the ratio of linear to
nonlinear work grows with width and is independent of depth.
\end{proposition}

\begin{proof}
The matrix $\mathbf{W}_\ell$ has $n_\ell n_{\ell-1}$ entries and each contributes
one multiply--accumulate; $\sigma$ is applied once per coordinate of each hidden
layer, of which there are $n_\ell$ at layer $\ell < L$. The stated asymptotics
follow by inspection: doubling every width multiplies the first sum by four and
the second by two, while appending a layer adds one term to each.
\end{proof}

Proposition~\ref{prop:nonlin-count} is the reason our measurements report both
quantities separately, and the reason the same network can look
nonlinearity-dominated or arithmetic-dominated depending on its shape: at small
width the $\sigma$ applications dominate the cost even though they are
asymptotically fewer, because each costs orders of magnitude more than a
multiply--accumulate under secret sharing.

\subsection{Nonlinearities, and which ones a cryptosystem can evaluate}
\label{sec:ml:activation}

\begin{definition}
\label{def:relu}
The \ib{rectified linear unit} is $\ReLU(x) = \max(x,0)$. The \ib{square}
activation is $\mathrm{sq}(x) = x^2$. A \ib{polynomial activation} of degree $d$
is any $p \in \Rr[x]$ with $\deg p = d$.
\end{definition}

\begin{proposition}
\label{prop:relu-props}
$\ReLU$ is $1$-Lipschitz and non-expansive: $\abs{\ReLU(x)-\ReLU(y)} \le
\abs{x-y}$ and $\abs{\ReLU(x)} \le \abs{x}$. Consequently, in a network whose
affine maps have operator norm at most $\gamma$ in the infinity norm, the
activations satisfy $\norm{\mathbf{u}_\ell}_\infty \le \gamma^{\ell}
\norm{\mathbf{u}_0}_\infty + O(\max_\ell \norm{\mathbf{b}_\ell}_\infty)$. By
contrast $\mathrm{sq}$ satisfies $\norm{\mathbf{u}_{\ell+1}}_\infty \le
(\gamma\norm{\mathbf{u}_\ell}_\infty)^2$, so for $\gamma\norm{\mathbf{u}}_\infty >
1$ the magnitude grows doubly exponentially in depth.
\end{proposition}

\begin{proof}
Both properties of $\ReLU$ follow from $\max(\cdot,0)$ being a projection onto a
convex set, or directly by cases. The propagation bound is then induction on
$\ell$: $\norm{\sigma(A_\ell \mathbf{u})}_\infty \le \norm{A_\ell
\mathbf{u}}_\infty \le \gamma \norm{\mathbf{u}}_\infty +
\norm{\mathbf{b}_\ell}_\infty$. For the square activation the same induction gives
$\norm{\mathbf{u}_{\ell+1}}_\infty = \norm{(A_\ell\mathbf{u}_\ell)^2}_\infty \le
(\gamma\norm{\mathbf{u}_\ell}_\infty)^2$, and iterating a map of the form
$m \mapsto (\gamma m)^2$ from $m_0$ gives $\log m_\ell = 2^{\ell}(\log m_0 + \log
\gamma) - \log\gamma$, which diverges when $\gamma m_0 > 1$ and converges to zero
when $\gamma m_0 < 1$.
\end{proof}

Proposition~\ref{prop:relu-props} is the analytic content behind a practical
constraint we measure. A fixed-point or lattice representation carries a bounded
dynamic range, so an activation whose iterates diverge doubly exponentially
exhausts that range within a few layers, and one whose iterates converge to zero
erases the signal. The square activation has a neutral point at
$\gamma\norm{\mathbf{u}}_\infty = 1$ and is unstable on both sides of it, whereas
$\ReLU$ is neutral everywhere. Standard initialisation, which sets
$\norm{\mathbf{W}}$ to approximately $1/\sqrt{n_{\ell-1}}$ so that the layer has
unit gain, places a square-activation network at precisely that neutral point,
which is why such networks are trainable at all and why they must be trained
rather than converted.

The cryptographic significance is a dichotomy.

\begin{remark}
\label{rem:activation-dichotomy}
Evaluating $\ReLU$ requires the \emph{sign} of a hidden value, that is, a
comparison. A secret-sharing protocol can compute a comparison exactly, at the
cost of Corollary~\ref{cor:cmp-cost}. An approximate homomorphic scheme over the
reals cannot compute a comparison at all: its operations are addition and
multiplication of ciphertexts, and the sign function is not a polynomial. Such a
scheme must therefore replace $\sigma$ by a polynomial approximation, paying in
multiplicative depth for the degree and in accuracy for the residual error, or the
network must be trained with a polynomial activation in the first place. The only
activations both paradigms evaluate \emph{exactly} are polynomials, and the
cheapest nonconstant one is $\mathrm{sq}$. This dichotomy, and not any difference
in raw throughput, determines the shape of our matched comparison in
Section~\ref{sec:matched}.
\end{remark}

\subsection{Structured layers: convolution and pooling}
\label{sec:ml:conv}

Image classifiers do not use dense affine maps at every layer. The structured
alternatives matter here because their structure is exactly what a packed
ciphertext can exploit.

\begin{definition}
\label{def:conv}
Let $\mathbf{u} \in \Rr^{c \times h \times w}$ be an input with $c$ channels and
let $\mathbf{K} \in \Rr^{c' \times c \times k \times k}$ be a kernel. The
\ib{convolution} with unit stride and no padding is
\[
  (\mathbf{K} * \mathbf{u})_{o,y,x} \;=\;
  \sum_{i=0}^{c-1}\ \sum_{a=0}^{k-1}\ \sum_{b=0}^{k-1}
  \mathbf{K}_{o,i,a,b}\, \mathbf{u}_{i,\,y+a,\,x+b} ,
\]
for $o < c'$ and $y < h-k+1$, $x < w-k+1$. \ib{Max pooling} with window $s$
replaces each disjoint $s \times s$ window by its maximum, and \ib{average
pooling} by its mean.
\end{definition}

\begin{proposition}
\label{prop:conv-linear}
Convolution is linear in $\mathbf{u}$, and for fixed $\mathbf{K}$ it is a
matrix--vector product with a matrix whose entries are the kernel entries,
repeated with a shift structure. It performs $c'c k^2 (h-k+1)(w-k+1)$
multiply--accumulate operations, and it applies no nonlinearity.
\end{proposition}

\begin{proof}
Linearity is immediate from the definition, which is a finite sum of scalar
multiples of coordinates of $\mathbf{u}$. Flattening $\mathbf{u}$ and the output
into vectors exhibits the map as a matrix; the count is the number of terms in the
triple sum times the number of output positions.
\end{proof}

Proposition~\ref{prop:conv-linear} says that a convolutional layer is free of
nonlinearity, so by Proposition~\ref{prop:nonlin-count} its cost under secret
sharing is traffic and not rounds. Its interaction with packing is sharper:
Equation~\eqref{eq:negconv} showed that a product in a negacyclic ring is a
shifted-and-signed convolution of coefficient vectors, and
Definition~\ref{def:conv} is a shifted inner product. Laying the input forwards
and the kernel reversed therefore places an entire convolutional layer in the
coefficients of a \emph{single} ring product, with the wanted outputs at a stride
and the unwanted ones discarded. That is the mechanism behind the technique we
implement and measure, and it explains why the ciphertext cost of a convolution in
our harness does not grow with the spatial size of the image.

Max pooling, by contrast, is a nonlinearity: it is a fold of maxima, and a maximum
is a mux on a comparison, $\max(x,y) = y + (x-y)\cdot[x>y]$. It therefore falls on
the expensive side of the dichotomy of Remark~\ref{rem:activation-dichotomy}, and
a window of size $s$ costs $s^2-1$ comparisons per output. Average pooling is
linear and free. This asymmetry is a design lever: our measurements exploit the
fact that $\ReLU$ and max pooling commute, so pooling first reduces the number of
$\ReLU$ applications by the pooling factor.

\begin{proposition}
\label{prop:pool-commute}
For any $s$, $\max$ over a window commutes with a coordinatewise nondecreasing
map. In particular $\mathrm{maxpool} \circ \ReLU = \ReLU \circ \mathrm{maxpool}$,
and more generally $\mathrm{maxpool} \circ (\ReLU \circ \rho) = \ReLU \circ \rho
\circ \mathrm{maxpool}$ for any nondecreasing rescaling $\rho$.
\end{proposition}

\begin{proof}
Let $f$ be nondecreasing and $S$ a finite set. Then $\max_{x \in S} f(x) =
f(\max_{x\in S} x)$: the right side is attained by some $x^\ast$ maximising $x$,
and $f(x^\ast) \ge f(x)$ for every $x \in S$ by monotonicity, so it is also the
maximum of the image; conversely it lies in the image. Both $\ReLU$ and a
nonnegative rescaling are nondecreasing, as is their composition.
\end{proof}

Proposition~\ref{prop:pool-commute} converts directly into a saving, because
applying the nonlinearity after pooling evaluates it on $s^2$ times fewer values.
Ordering the two operations is therefore not a matter of taste; it changes the
comparison count of the whole network, and we report the measured factor.

\subsection{The decision: why classification does not need a softmax}
\label{sec:ml:argmax}

Trained classifiers are usually described as ending in a \ib{softmax},
$\mathrm{sm}(\mathbf{z})_j = e^{z_j}/\sum_l e^{z_l}$, which converts logits into a
probability vector. Under encryption this would be a serious obstacle: it needs an
exponential and a reciprocal, neither of which is a low-degree polynomial. It is
also unnecessary.

\begin{proposition}
\label{prop:argmax-invariance}
Softmax is strictly increasing in each coordinate relative to the others, in the
sense that $z_j > z_l \iff \mathrm{sm}(\mathbf{z})_j > \mathrm{sm}(\mathbf{z})_l$.
Consequently $\argmax \circ\, \mathrm{sm} = \argmax$, and a classifier that
reveals only its predicted label may omit the softmax entirely. The same holds for
any strictly increasing scalar function applied coordinatewise, and for
multiplication by any positive constant.
\end{proposition}

\begin{proof}
The denominator of $\mathrm{sm}$ is a positive constant common to all coordinates,
and $t \mapsto e^{t}$ is strictly increasing, so
$\mathrm{sm}(\mathbf{z})_j > \mathrm{sm}(\mathbf{z})_l$ iff $e^{z_j} > e^{z_l}$ iff
$z_j > z_l$. Hence the coordinatewise order, and in particular its maximiser, is
preserved. For a general strictly increasing $f$ applied coordinatewise the same
argument applies verbatim; for a positive scalar multiple it is immediate.
\end{proof}

Proposition~\ref{prop:argmax-invariance} is what makes a \emph{classifier}
tractable under encryption while a language model is not: the former reveals
$\argmax$ of the logits, which is invariant under the transformation one would
otherwise have to compute, whereas the latter must sample from the distribution
and therefore needs the normalisation itself. It also fixes what a private
inference protocol should output, namely a sharing or ciphertext of the index, and
what it should reveal, namely that index alone.

\begin{definition}
\label{def:argmax-prot}
Given logits $\mathbf{z} \in \Rr^{C}$, the \ib{argmax fold} computes
$(m_0,i_0) = (z_0,0)$ and
\[
  c_j = [\,z_j > m_{j-1}\,], \qquad
  m_j = m_{j-1} + c_j (z_j - m_{j-1}), \qquad
  i_j = i_{j-1} + c_j (j - i_{j-1}),
\]
returning $i_{C-1}$. It uses $C-1$ comparisons and $2(C-1)$ multiplications, and
its comparisons are sequentially dependent.
\end{definition}

The sequential dependence in Definition~\ref{def:argmax-prot} is worth noting
because it is the one place in an inference where our round count is not
determined by network depth: a $C$-way argmax contributes $C-1$ dependent
comparisons. A tournament fold reduces this to $\ceil{\log C}$ levels at the same
comparison count, which matters when $C$ is large.

\subsection{Quantisation: from real networks to the arithmetic a protocol offers}
\label{sec:ml:quant}

Networks are trained over the reals; protocols offer $\Z_{2^k}$, $\mathbb{F}_p$ or
approximate reals. The bridge is quantisation, and because its error is one of the
two error terms in every accuracy figure we report, we state it explicitly.

\begin{definition}
\label{def:qnet}
Fix a scale $f$. The \ib{quantised network} $\widehat{g}$ of a network $g$
replaces every parameter by its fixed-point encoding at scale $f$
(Definition~\ref{def:fixed}) and every affine map by integer arithmetic followed
by a rescale, so that layer $\ell$ computes
$\widehat{\mathbf{u}}_\ell = \sigma\bigl(\rnd{2^{-f}\widehat{\mathbf{W}}_\ell
\widehat{\mathbf{u}}_{\ell-1}} + \widehat{\mathbf{b}}_\ell\bigr)$, with the bias
stored at output scale.
\end{definition}

\begin{proposition}
\label{prop:quant-error}
Let $g$ have $L$ layers of width at most $n$, with parameters and activations
bounded by $M$ and each $\sigma$ $1$-Lipschitz. Then the quantised network
satisfies
\[
  \norm{\widehat{g}(x)\,2^{-f} - g(x)}_\infty \;\le\;
  \bigl(nM + \tfrac12\bigr)\,\frac{\gamma^{L}-1}{\gamma-1}\,2^{-f}
  \qquad\text{where } \gamma = nM,
\]
that is, the per-layer error of Proposition~\ref{prop:fixed-error} amplified by
the network's gain and accumulated over depth.
\end{proposition}

\begin{proof}
By Proposition~\ref{prop:fixed-error} each layer introduces error at most
$\varepsilon = (nM + \tfrac12)2^{-f}$ relative to exact evaluation on its own
input. Write $\eta_\ell$ for the error at the output of layer $\ell$. A
$1$-Lipschitz $\sigma$ does not amplify, and the affine map amplifies by at most
its infinity-norm operator norm, bounded by $nM$; hence $\eta_\ell \le \gamma
\eta_{\ell-1} + \varepsilon$ with $\eta_0 = 0$. Unrolling the recurrence gives
$\eta_L \le \varepsilon(1 + \gamma + \cdots + \gamma^{L-1})$, which is the stated
geometric sum.
\end{proof}

Proposition~\ref{prop:quant-error} is pessimistic---it assumes worst-case sign
alignment at every step, whereas real errors partly cancel---and we therefore do
not use it to predict accuracy. We use it for its shape, which carries two
practical consequences. Error grows geometrically in depth with ratio equal to the
network's gain, so a network with unit gain accumulates error linearly in depth
while one with gain above unity does not remain usable; and the scale $f$ trades
directly against every cost in the protocol, because the arithmetic domain must
carry $f$ plus the integer range plus the statistical margin of
Remark~\ref{rem:kappa-cost}. In practice one measures not bounds, which is
why our pipeline reports the observed quantisation error of every model it
deploys.

\begin{remark}
\label{rem:exactness}
There is a stronger property than small error, and our harness insists on it
wherever it is available: \ib{bit-exactness}. A private evaluation is bit-exact
when its output equals that of the \emph{same} quantised network evaluated in the
clear, digit for digit. Bit-exactness says nothing about whether quantisation was
a good idea; it says that the cryptography computed what the plaintext reference
computed, and it therefore separates a protocol error from an approximation error.
Where a paradigm cannot offer it---an approximate homomorphic scheme cannot---the
weaker comparison is agreement of the predicted label, and we say which of the two
a figure rests on.
\end{remark}

\subsection{What an accuracy figure means}
\label{sec:ml:metrics}

Two of our case studies are medical, where accuracy alone is misleading, so we fix
the metrics.

\begin{definition}
\label{def:metrics}
For a binary classifier with true positives $\mathit{TP}$, false positives
$\mathit{FP}$, true negatives $\mathit{TN}$ and false negatives $\mathit{FN}$,
\ib{accuracy} is $(\mathit{TP}+\mathit{TN})/N$, \ib{sensitivity} (recall) is
$\mathit{TP}/(\mathit{TP}+\mathit{FN})$, \ib{specificity} is
$\mathit{TN}/(\mathit{TN}+\mathit{FP})$, and \ib{balanced accuracy} is the mean of
the last two. For $C$ classes, balanced accuracy is the mean of the per-class
recalls. The \ib{area under the receiver operating characteristic}, $\mathrm{AUC}$,
is the probability that a uniformly chosen positive receives a higher score than a
uniformly chosen negative, ties counting one half.
\end{definition}

\begin{proposition}
\label{prop:auc-mw}
$\mathrm{AUC}$ equals the normalised Mann--Whitney statistic: if $s_1,\dots,s_m$
are the scores of the positives and $t_1,\dots,t_n$ those of the negatives, then
\[
  \mathrm{AUC} \;=\; \frac{1}{mn}\sum_{i=1}^{m}\sum_{j=1}^{n}
  \Bigl(\mathbb{1}[s_i > t_j] + \tfrac12\,\mathbb{1}[s_i = t_j]\Bigr).
\]
\end{proposition}

\begin{proof}
The double sum counts, over all $mn$ pairs, the indicator of the event in
Definition~\ref{def:metrics} with ties weighted one half; dividing by the number
of pairs converts the count into the probability under the uniform choice of a
pair.
\end{proof}

Proposition~\ref{prop:auc-mw} is the form an implementation must use, and the tie
term is the part that is easy to omit; a quantised network produces exact ties far
more often than a real-valued one, because its scores are integers, so an
$\mathrm{AUC}$ computed without the tie correction is biased for exactly the models
we deploy.

\begin{definition}
\label{def:wilson}
For $k$ successes in $N$ independent trials with unknown rate $p$, the
\ib{Wilson score interval} at level $1-\alpha$ is
\[
  \frac{1}{1+z^2/N}\left(\widehat{p} + \frac{z^2}{2N}
  \;\pm\; z\sqrt{\frac{\widehat{p}(1-\widehat{p})}{N} + \frac{z^2}{4N^2}}\right),
  \qquad \widehat{p} = k/N,
\]
where $z$ is the standard normal quantile at $1-\alpha/2$.
\end{definition}

We quote Wilson instead of normal-approximation intervals because our evaluation
sets are small enough, and our rates close enough to one, that the normal interval
can extend past unity; Wilson does not, and it retains nominal coverage at those
extremes. An accuracy difference between two configurations is reported as
meaningful only when the intervals are disjoint.

\subsection{Threat models for private inference}
\label{sec:ml:threat}

Finally, the same computation admits several confidentiality goals, and the
literature's figures are not comparable across them. We name them because our
comparison depends on holding one fixed.

\begin{definition}
\label{def:pi-models}
In \ib{input privacy} the client's $x$ is hidden from the server and the
parameters $\theta$ are public or held in the clear by the server. In \ib{model
privacy} $\theta$ is hidden from the client. In \ib{joint privacy} both are
hidden from whoever evaluates: the evaluator learns neither $x$ nor $\theta$, and
only the label is revealed, to a designated recipient.
\end{definition}

\begin{remark}
\label{rem:pi-cost}
The distinction is not a detail of framing; it changes the arithmetic. Under input
privacy alone, an affine layer is a \emph{plaintext}--ciphertext product, which for
a lattice scheme costs no ciphertext multiplication and adds far less noise
(Proposition~\ref{prop:noise-growth}(2)) and for a secret sharing is a purely
local scaling. Under joint privacy it is a ciphertext--ciphertext or
sharing--sharing product, which consumes multiplicative depth or a Beaver triple
per term. Comparing a joint-privacy system against an input-privacy one on
latency therefore compares two different computations, and the hybrid systems most
often cited for their speed---in which the server holds the model in the
clear---occupy the cheaper setting. Our measurements are of joint privacy
throughout, and where we place them beside a published figure we state which
setting that figure was obtained in.
\end{remark}
